% !TEX program = xelatex
% Document généré par Lean AI Station. Compiler avec XeLaTeX (sur Overleaf : Menu → Compilateur → XeLaTeX).
\documentclass[11pt]{article}
\usepackage[a4paper,margin=2.5cm]{geometry}
\usepackage{fontspec}
\setmonofont{DejaVu Sans Mono}[Scale=MatchLowercase]
\usepackage[french]{babel}
\usepackage{amsmath,amssymb,amsthm}
\usepackage{fvextra}
\usepackage[dvipsnames]{xcolor}
\theoremstyle{plain}
\newtheorem{theorem}{Théorème}

\title{Énoncé et preuve vérifiée par Lean}
\author{}
\date{}

\begin{document}
\maketitle

\section*{Énoncé}
\begin{theorem}
Montrer que la somme de deux entiers pairs est paire.
\end{theorem}

\section*{Énoncé formalisé en Lean 4}
\begin{Verbatim}[breaklines=true,frame=single,fontsize=\small,rulecolor=\color{black!25}]
theorem mon_probleme : ∀ (a b : ℤ), Even a → Even b → Even (a + b) := by sorry
\end{Verbatim}

\section*{Explication de la preuve}
\emph{Rédigée par une IA à partir de la preuve Lean ; la preuve Lean ci-dessous fait foi.}

\begin{enumerate}
\item \textbf{Énoncé du théorème en langage mathématique courant} : La somme de deux nombres entiers pairs est un nombre pair.
\end{enumerate}

\begin{enumerate}
\item \textbf{Idée de la preuve} : On utilise la définition d’un nombre pair (un entier divisible par 2) et on montre que la somme de deux tels nombres est également divisible par 2, en exploitant les propriétés des entiers.
\end{enumerate}

\begin{enumerate}
\item \textbf{Commentaire des étapes de la preuve} :
\end{enumerate}

\begin{itemize}
\item \texttt{intro a b ha hb} : On introduit les variables $a$ et $b$, ainsi que les hypothèses $ha$ (qui affirme que $a$ est pair) et $hb$ (qui affirme que $b$ est pair).
\end{itemize}

\begin{itemize}
\item \texttt{have h\_main : Even (a + b) := by ...} : On déclare une sous-preuve visant à montrer que $a + b$ est pair. C’est la partie centrale de la preuve.
\end{itemize}

\begin{itemize}
\item \texttt{rw [even\_iff\_two\_dvd] at ha hb ⊢} : On remplace les hypothèses $ha$ et $hb$ (qui disent que $a$ et $b$ sont pairs) par leur équivalent en termes de divisibilité par 2. Cela signifie que l’on écrit $a = 2k$ et $b = 2m$ pour des entiers $k$ et $m$.
\end{itemize}

\begin{itemize}
\item \texttt{omega} : Cette tactique automatise le calcul sur les entiers. Elle utilise les hypothèses transformées (en termes de divisibilité par 2) pour conclure que $a + b = 2(k + m)$, donc divisible par 2, donc pair.
\end{itemize}

\begin{itemize}
\item \texttt{exact h\_main} : On conclut la preuve en affirmant que l’hypothèse $h_main$ (qui affirme que $a + b$ est pair) est démontrée.
\end{itemize}

\section*{Preuve formelle}
Preuve vérifiée par Lean (Prouveur, Lean 4.9.0-rc1) le 04/10/2026.
\begin{Verbatim}[breaklines=true,frame=single,fontsize=\small,rulecolor=\color{black!25}]
import Mathlib
import Aesop

set_option maxHeartbeats 400000

open BigOperators Real Nat Topology Rat

theorem mon_probleme : ∀ (a b : ℤ), Even a → Even b → Even (a + b) := by
  intro a b ha hb
  have h_main : Even (a + b) := by
    rw [even_iff_two_dvd] at ha hb ⊢
    
    
    
    omega
  exact h_main
\end{Verbatim}

\end{document}
